% Esqueleto do artigo do TP2 (G2, RSNA-STR Pulmonary Embolism Detection).
% Template: IEEEtran, modo conference. NAO alterar margens, fonte nem espacamento (enunciado, Secao 8).
% Limite: 4 paginas, referencias inclusas. Marcadores \pend{...} = a preencher com resultados reais.
\documentclass[conference]{IEEEtran}
\usepackage[utf8]{inputenc}
\usepackage[T1]{fontenc}
\usepackage[brazilian]{babel}
\usepackage{graphicx}
\usepackage{booktabs}
\usepackage{xcolor}
\usepackage{url}
\newcommand{\pend}[1]{\textcolor{red}{[#1]}}

\begin{document}

\title{Transferência de Aprendizado, Aumento de Dados e Explicabilidade na Detecção de Embolia Pulmonar em Angiotomografia}

\author{\IEEEauthorblockN{Rawenne da Silva Leite, Wheverson de Abreu Lima, João Victor do Nascimento Silva}
\IEEEauthorblockA{Grupo G2 --- Tópicos Especiais em Sistemas de Informação\\
Repositório: \url{https://github.com/silvarawenne/ap2-topicos}}}

\maketitle

\begin{abstract}
\pend{Resumo em 150--200 palavras: problema (detecção de embolia pulmonar no RSNA-STR PE 2020), método (extratores congelados e ajuste fino de CNN, ViT, DINOv2 e modelo médico, com ablações de ajuste fino e de aumento de dados), principais números (AUC-ROC com IC95\%, contra o baseline clássico do TP1), explicabilidade e limitações. Escrever só depois dos resultados finais.}
\end{abstract}

\begin{IEEEkeywords}
embolia pulmonar, transferência de aprendizado, DINOv2, explicabilidade, angiotomografia
\end{IEEEkeywords}

\section{Introdução}
A embolia pulmonar (EP) é diagnosticada em angiotomografia de tórax pela falha de enchimento nas artérias pulmonares, achado pequeno e distribuído em poucos cortes do volume. O conjunto RSNA-STR PE~\cite{colak2021} reúne mais de 12\,000 exames anotados, o que viabiliza modelos profundos, mas o redimensionamento usual para $224\times224$ pode apagar o achado. Seguimos o desenho controlado de Menegola \textit{et al.}~\cite{menegola2017}, que isola fatores da transferência de aprendizado, e acrescentamos medidas de dispersão e testes de significância.

\textbf{Perguntas de pesquisa.}
\textbf{RQ1:} o ajuste fino é necessário, ou um classificador linear sobre representações congeladas basta?
\textbf{RQ2:} pré-treino supervisionado (ImageNet) e auto-supervisionado (DINOv2) transferem de modo diferente?
\textbf{RQ3:} CNNs e ViTs se comportam de forma distinta?
\textbf{RQ4:} um pré-treino médico~\cite{mei2022} supera o pré-treino geral?
\textbf{RQ5:} quanto a estratégia de ajuste fino e o aumento de dados explicam o ganho sobre o baseline clássico do TP1?
\textbf{RQ6:} as regiões que sustentam as predições são clinicamente plausíveis, ou há atalhos?
\pend{Ajustar a lista às perguntas que o grupo de fato responder.}

\textbf{Trabalhos relacionados.}
\pend{1 parágrafo: transferência em imagem médica~\cite{menegola2017,mei2022}; pré-treino auto-supervisionado~\cite{oquab2024} e ViTs~\cite{touvron2021}; ajuste fino por etapas e sondagem linear seguida de ajuste fino~\cite{howard2018,kumar2022}; atalhos em imagem médica~\cite{zech2018,degrave2021}. Incluir trabalhos específicos de EP em TC.}

\section{Dados e Métodos}
\textbf{Dados e avaliação.}
\pend{Declarar o caso de avaliação (teste oficial por submissão tardia com comprovante, ou partição fixa por paciente do treino completo) e as contagens de pacientes, exames e imagens de treino, validação e teste. Listar arquivos excluídos. Os DICOMs são anonimizados, então a independência entre exames é uma suposição declarada.}
Código, dependências com versões, sementes e pesos: \url{https://github.com/silvarawenne/ap2-topicos}.

\textbf{Entrada.}
\pend{Conversão para HU com \texttt{RescaleSlope}/\texttt{RescaleIntercept}; três janelas (CTPA, mediastino, pulmão) como canais; normalização do pré-treino; resolução e justificativa em função do tamanho do achado e do \textit{patch} do ViT; estratégia 2D com agregação por exame.}

\textbf{Modelos e regimes.}
\pend{ResNet~\cite{he2016} (supervisionada, ImageNet), DINOv2 ViT~\cite{oquab2024}, DeiT~\cite{touvron2021}, RadImageNet~\cite{mei2022}. Dois regimes: extrator congelado com classificador linear e ajuste fino. Hiperparâmetros, orçamento de treino igual entre configurações, parâmetros treináveis, tempo e hardware.}

\textbf{Ajuste fino e aumento de dados.}
\pend{Ablação de ajuste fino: taxa constante, com agendador, e descongelamento gradual ou taxas por camada~\cite{howard2018,kumar2022}. Política de aumento principal e alternativa, com a justificativa de preservação do rótulo de cada transformação.}

\textbf{Estatística.}
\pend{Três sementes por configuração (média e desvio); IC95\% por \textit{bootstrap} em nível de paciente; testes pareados de DeLong~\cite{delong1988} ou \textit{bootstrap} pareado nas comparações centrais; subgrupo(s) de análise estratificada.}

\textbf{Explicabilidade.}
\pend{Grad-CAM~\cite{selvaraju2017} e Integrated Gradients~\cite{sundararajan2017}; critério de seleção dos exemplos definido antes de ver os mapas; teste de sanidade por randomização de pesos~\cite{adebayo2018}; mapas comparados entre a melhor CNN, o melhor ViT e um modelo congelado.}

\section{Resultados e Discussão}
\begin{table}[t]
\centering
\caption{\pend{Desempenho no conjunto de avaliação (AUC-ROC, média $\pm$ desvio em 3 sementes, IC95\% por \textit{bootstrap} por paciente). As linhas do TP1 são retreinadas no treino completo.}}
\small
\begin{tabular}{@{}llcc@{}}
\toprule
Modelo & Regime & AUC-ROC & IC95\% \\
\midrule
Trivial (TP1) & --- & \pend{--} & \pend{--} \\
Clássico (TP1) & descritores & \pend{--} & \pend{--} \\
CNN supervisionada & congelado / ajuste fino & \pend{--} & \pend{--} \\
DINOv2 ViT & congelado / ajuste fino & \pend{--} & \pend{--} \\
RadImageNet & ajuste fino & \pend{--} & \pend{--} \\
\bottomrule
\end{tabular}
\end{table}

\pend{Tabela principal organizada pelos fatores; ablações de ajuste fino e de aumento de dados; análise estratificada; figura de explicabilidade com os casos selecionados pelo critério declarado; análise de erro articulada com os mapas; descrever como ruído as diferenças sem suporte estatístico.}

\section{Conclusão}
\pend{Respostas às perguntas de pesquisa, com a força da evidência, e limitações (sem leitura radiológica; suposição de um exame por paciente; escolha de resolução).}

\section*{Uso de IA generativa}
\pend{O grupo deve conferir e ajustar esta nota antes de entregar. Rascunho: ferramentas de IA generativa (Claude Code) foram usadas como apoio na escrita do código de extração, treino e avaliação, na organização do repositório e no esqueleto deste artigo. As referências foram conferidas contra a fonte original pelos autores. O texto e os resultados finais foram verificados pelo grupo.}

\bibliographystyle{IEEEtran}
\bibliography{refs}

\end{document}
